\documentclass[aspectratio=169,11pt]{beamer}
% Shared Beamer style for practice contest solution walkthroughs.
\usepackage[utf8]{inputenc}
\usepackage[T1]{fontenc}
\usepackage{lmodern}
\usepackage{amsmath,amssymb}
\usepackage{xcolor}
\usepackage{hyperref}
\usepackage{listings}

\definecolor{CPNavy}{HTML}{172033}
\definecolor{CPBlue}{HTML}{2563EB}
\definecolor{CPGreen}{HTML}{16A34A}
\definecolor{CPOrange}{HTML}{F97316}
\definecolor{CPGray}{HTML}{64748B}
\definecolor{CPLight}{HTML}{F8FAFC}
\definecolor{CPCodeBg}{HTML}{F1F5F9}
\definecolor{CPCodeRule}{HTML}{CBD5E1}

\usetheme{default}
\usefonttheme{professionalfonts}
\setbeamertemplate{navigation symbols}{}
\setbeamersize{text margin left=0.65cm,text margin right=0.65cm}
\setbeamertemplate{itemize item}{\color{CPBlue}$\blacktriangleright$}
\setbeamertemplate{itemize subitem}{\color{CPGray}$\triangleright$}

\setbeamercolor{normal text}{fg=CPNavy,bg=white}
\setbeamercolor{frametitle}{fg=white,bg=CPNavy}
\setbeamercolor{title}{fg=white,bg=CPNavy}
\setbeamercolor{subtitle}{fg=CPLight,bg=CPNavy}
\hypersetup{colorlinks=true,urlcolor=CPBlue}

\setbeamertemplate{frametitle}{%
  \nointerlineskip%
  \begin{beamercolorbox}[wd=\paperwidth,ht=0.88cm,dp=0.22cm,leftskip=0.55cm,rightskip=0.35cm]{frametitle}
    \usebeamerfont{frametitle}\insertframetitle\hfill\small\insertframenumber/\inserttotalframenumber
  \end{beamercolorbox}%
}

\setbeamertemplate{footline}{%
  \leavevmode%
  \hbox{%
    \begin{beamercolorbox}[wd=.58\paperwidth,ht=2.4ex,dp=1ex,leftskip=.5cm]{author in head/foot}%
      \color{CPGray}\insertshorttitle
    \end{beamercolorbox}%
    \begin{beamercolorbox}[wd=.42\paperwidth,ht=2.4ex,dp=1ex,rightskip=.5cm plus1fil]{date in head/foot}%
      \hfill\color{CPGray}\insertshortauthor
    \end{beamercolorbox}%
  }%
}

\setbeamertemplate{title page}{%
  \vspace*{1.25cm}
  \begin{beamercolorbox}[wd=\paperwidth,sep=0.65cm,leftskip=0.15cm,rightskip=0.15cm]{title}
    {\color{white}\Huge\bfseries\inserttitle\par}
    \vspace{0.35cm}
    {\color{CPLight}\Large\insertsubtitle\par}
    \vspace{0.6cm}
    {\color{CPLight}\large\insertauthor\par}
  \end{beamercolorbox}
  \vfill
}

\newcommand{\sectiontag}[1]{\textbf{\color{CPBlue}#1}}
\newenvironment{tightitemize}{\begin{itemize}\small\setlength{\itemsep}{0.18cm}}{\end{itemize}}
\lstdefinestyle{cpcode}{
  language=Python,
  basicstyle=\ttfamily\tiny,
  keywordstyle=\color{CPBlue}\bfseries,
  commentstyle=\color{CPGray}\itshape,
  stringstyle=\color{CPGreen},
  backgroundcolor=\color{CPCodeBg},
  frame=single,
  framerule=0.25pt,
  rulecolor=\color{CPCodeRule},
  showstringspaces=false,
  columns=fullflexible,
  keepspaces=true,
  breaklines=true,
  breakatwhitespace=false,
  tabsize=4,
  xleftmargin=0.08cm,
  xrightmargin=0.08cm,
  aboveskip=0.08cm,
  belowskip=0.08cm
}


\title[pairingsocks]{Pairing Socks}
\subtitle{Day 1 solution walkthrough}
\author[Solution Deck]{Competitive Programming Bootcamp}
\date{}

\begin{document}

\begin{frame}[plain]
  \titlepage
\end{frame}

% BEGIN GENERATED STATEMENT INTERPRETATION
\begin{frame}{Interpret the Word Problem}
\small
\sectiontag{Statement excerpts:}
\begin{tightitemize}
  \item \textit{``socks stacked in a pile''}
  \item \textit{``least number of moves she needs to pair the socks''}
\end{tightitemize}

\vspace{0.18cm}
\sectiontag{Programming interpretation:}
\begin{tightitemize}
  \item Model the original pile and auxiliary pile as stacks.
  \item A move transfers one top sock; a pair operation removes matching top socks.
  \item If the deterministic forward simulation empties both stacks, the move count is the answer.
\end{tightitemize}
\end{frame}
% END GENERATED STATEMENT INTERPRETATION







\begin{frame}{Problem Model}
\small
\sectiontag{Textbook connection:} CP4 2.2 j. Stack

\vspace{0.25cm}
\sectiontag{Core technique:} Stack simulation

\vspace{0.25cm}
\begin{tightitemize}
  \item The input order is a stack of socks.
  \item A second stack represents the auxiliary pile.
  \item A pair disappears when the top sock of each stack has the same color.
\end{tightitemize}
\end{frame}

\begin{frame}{How to Recognize the Approach}
\begin{tightitemize}
  \item Only the top of each pile matters, so stack operations are natural.
  \item Moving a sock back from auxiliary to original does not create a useful new order.
  \item The process is deterministic once we always move from original to auxiliary.
\end{tightitemize}
\end{frame}

\begin{frame}{Algorithm}
\begin{tightitemize}
  \item Read the original stack and initialize an empty auxiliary stack.
  \item Pop one sock from the original stack and push it onto the auxiliary stack.
  \item After each move, repeatedly remove matching top socks from both stacks.
  \item Count every move or successful pair-removal action as required by the problem.
  \item If both stacks are empty, print the count; otherwise print impossible.
\end{tightitemize}
\end{frame}

\begin{frame}{Walkthrough of the Key State}
\begin{tightitemize}
  \item The outer loop keeps feeding the auxiliary stack.
  \item The inner loop handles any chain reaction of exposed matching socks.
  \item A nonempty auxiliary stack at the end means some socks could not be paired legally.
\end{tightitemize}
\end{frame}

\begin{frame}{Why the Algorithm Is Correct}
\begin{tightitemize}
  \item The algorithm simulates the only useful direction of movement.
  \item Whenever a pair is exposed, removing it immediately cannot block a future match because it only reveals deeper socks.
  \item If the process ends with leftover socks, no later legal action can pair them because the original stack is empty.
\end{tightitemize}
\end{frame}

\begin{frame}{Complexity and Implementation Notes}
\small
\sectiontag{Complexity:} O(N) time and O(N) memory.

\vspace{0.3cm}
\begin{tightitemize}
  \item Python lists are sufficient stacks with append and pop.
  \item Use truthiness of lists to test whether a stack is empty.
  \item Be careful that input lists are read in the orientation expected by the simulation.
\end{tightitemize}
\end{frame}



% BEGIN GENERATED CODE WALKTHROUGH
\begin{frame}[fragile]{Code: Initialize Two Stacks}
\scriptsize
\sectiontag{Source lines:} \texttt{pairingsocks.py:43--54}

\vspace{0.08cm}
\begin{columns}[T,totalwidth=\textwidth]
\begin{column}{0.58\textwidth}
\begin{lstlisting}[style=cpcode]
import sys

lines = sys.stdin.read().splitlines()
stack1 = lines[1].split()

sockPairs = int(lines[0])
moveCounter = 0
stack2 = []
\end{lstlisting}
\end{column}
\begin{column}{0.38\textwidth}
\sectiontag{What this chunk does}
\begin{tightitemize}
  \item The first line gives the number of pairs; the second line gives the stack order.
  \item stack1 is the original pile and stack2 starts empty.
  \item moveCounter tracks both transfers and pair removals.
\end{tightitemize}
\end{column}
\end{columns}
\end{frame}

\begin{frame}[fragile]{Code: Simulate Legal Moves}
\scriptsize
\sectiontag{Source lines:} \texttt{pairingsocks.py:56--63}

\vspace{0.08cm}
\begin{columns}[T,totalwidth=\textwidth]
\begin{column}{0.58\textwidth}
\begin{lstlisting}[style=cpcode]
while stack1:
    stack2.append(stack1.pop())
    moveCounter += 1

    while stack1 and stack2 and stack2[-1] == stack1[-1]:
        stack2.pop()
        stack1.pop()
        moveCounter += 1
\end{lstlisting}
\end{column}
\begin{column}{0.38\textwidth}
\sectiontag{What this chunk does}
\begin{tightitemize}
  \item Each loop moves one sock from the original pile to the auxiliary pile.
  \item Whenever the two top socks match, both are removed immediately.
  \item The inner loop handles chain reactions after exposing deeper socks.
\end{tightitemize}
\end{column}
\end{columns}
\end{frame}

\begin{frame}[fragile]{Code: Report Success or Failure}
\scriptsize
\sectiontag{Source lines:} \texttt{pairingsocks.py:65--68}

\vspace{0.08cm}
\begin{columns}[T,totalwidth=\textwidth]
\begin{column}{0.58\textwidth}
\begin{lstlisting}[style=cpcode]
if stack1 or stack2:
    sys.stdout.write("impossible")
else:
    sys.stdout.write(str(moveCounter))
\end{lstlisting}
\end{column}
\begin{column}{0.38\textwidth}
\sectiontag{What this chunk does}
\begin{tightitemize}
  \item Any remaining sock means a complete legal pairing was not found.
  \item If both stacks are empty, every sock was paired and the counted moves are printed.
\end{tightitemize}
\end{column}
\end{columns}
\end{frame}
% END GENERATED CODE WALKTHROUGH

\begin{frame}{Source}
\small
\sectiontag{Solution file:} \texttt{practice\_contests/day\_1/pairingsocks.py}

\vspace{0.3cm}
\sectiontag{Problem:} \url{https://open.kattis.com/problems/pairingsocks}

\vspace{0.45cm}
Use this deck with the source code open beside it. The slides explain the reasoning; the code shows the exact input handling and output formatting.
\end{frame}

\end{document}
